\documentclass{article}
\usepackage{graphicx}
\usepackage{amsmath}
\usepackage{amsfonts}
\usepackage{bookmark}
\usepackage{booktabs}
\usepackage{exercises}
\graphicspath{ {} }
\usepackage{indentfirst}
\usepackage{amsthm}
\usepackage{yhmath}
\usepackage{geometry}
\usepackage{placeins}
\usepackage{csquotes}
\usepackage{pgfplots}
\newcommand*{\oldneg}{\mathord{\sim}}
\newcommand\showdiv[1]{\overline{\smash{)}#1}}
\newcommand\Mydiv[1]{%
\smash{\raise.3ex\hbox{$\big)$}}\mkern-8mu
        \overline{\enspace\strut#1}}
\setlength{\tabcolsep}{5pt}
\theoremstyle{definition}
\newtheorem{definition}{Definition}

\setlength{\parindent}{.5in}
\setlength{\parskip}{12pt}
\renewcommand{\arraystretch}{2}
\geometry{margin=1in}
\newcommand{\?}{\stackrel{?}{=}}
\pgfplotsset{width=7cm,compat=1.18}\usepgfplotslibrary{statistics}

\title{9.4 Answers}
\begin{document}
\maketitle

\section*{1}
The inverse sine of $\frac{1}{2}$ is $\frac{\pi}{6}$, or $30^\circ$.

\section*{2}
$\tan\theta=\frac{opp}{adj}$ and $\cos\theta=\frac{adj}{hyp}$. Now, $hyp>adj$. But this means that 
$\frac{adj}{hyp}\ne 3$, and it is possible that $\frac{opp}{adj}=3$. Thus $\tan^{-1}3$ is defined, but 
$\cos^{-1}3$ is undefined.

\section*{3}
When $-90^\circ\le\theta\le 90^\circ$, or $-\frac{\pi}{2}\le\theta\le\frac{\pi}{2}$,
\[ \theta=\sin^{-1}1=90^\circ \] or $\frac{\pi}{2}$.

\section*{4}
When $-90^\circ\le\theta\le 90^\circ$, or $-\frac{\pi}{2}\le\theta\le\frac{\pi}{2}$,
\[ \theta=\tan^{-1}\left(-1\right)=-45^\circ \] or $\frac{\pi}{4}$.

\section*{5}
When $0^\circ\le\theta\le 180^\circ$, or $0\le\theta\le\pi$,
\[ \theta=\cos^{-1}0=90^\circ \] or $\frac{\pi}{2}$.

\section*{6}
$\cos^{-1}(-2)$ is undefined.

\section*{7}
When $-90^\circ\le\theta\le 90^\circ$, or $-\frac{\pi}{2}\le\theta\le\frac{\pi}{2}$, \[ \theta=\sin^{-1}\frac{\sqrt{3}}{2}=60^\circ \] or 
$\frac{\pi}{3}$.

\section*{8}
When $-90^\circ\le\theta\le 90^\circ$, or $-\frac{\pi}{2}\le\theta\le\frac{\pi}{2}$,
\[ \theta=\sin^{-1}\frac{1}{2}=30^\circ \] or $\frac{\pi}{6}$.

\section*{9}
When $-90^\circ\le\theta\le 90^\circ$, or $-\frac{\pi}{2}\le\theta\le\frac{\pi}{2}$,
\[ \theta=\tan^{-1}\left(-\frac{\sqrt{3}}{3}\right)=\tan^{-1}\left(-\frac{1}{\sqrt{3}}\right)=-30^\circ \] or
$-\frac{\pi}{6}$.

\section*{10}
When $0^\circ\le\theta\le 180^\circ$, or $0\le\theta\le\pi$,
\[ \theta=\cos^{-1}(\frac{1}{2})=60^\circ \] or $\frac{\pi}{3}$.

\section*{11}
When $0^\circ\le\theta\le 180^\circ$, or $0\le\theta\le\pi$,
\[ \theta=\cos^{-1}\frac{\sqrt{2}}{2}=\cos^{-1}\frac{1}{\sqrt{2}}=45^\circ \] or $\frac{\pi}{4}$.

\section*{16}
$\tan^{-1}(-0.75)\approx -36.87^\circ$ or $-0.64$.

\section*{17}
$\sin^{-1}(-0.2)\approx -11.54^\circ$ or $-0.20$.

\section*{19}
$\cos^{-1}0.99\approx 8.11^\circ$ or $0.14$.

\section*{22}
In the range of $-90^\circ\le\theta\le 90^\circ$,
\[ \theta=\sin^{-1} 0.15\approx 8.63^\circ \] An angle in the range $90^\circ <\theta<180^\circ$ is about
$180-8.63=171.37$.

\section*{24}
In the range of $-90^\circ\le\theta\le 90^\circ$,
\[ \theta=\tan^{-1}(-5.3)\approx -79.32 \] An angle in the range $90^\circ <\theta<180^\circ$ is about
$180-79.32=100.68$.

\section*{25}
In the range of $0^\circ\le\theta\le 180^\circ$,
\[ \theta=\cos^{-1}0.25\approx 75.52^\circ \] An angle in the range $270^\circ<\theta<360^\circ$ is about
$360-75.52=284.48^\circ$.

\section*{26}
The answer given is $\theta\approx 44.4^\circ$. But the question asked for an angle in the range $90^\circ<\theta<180^\circ$, meaning the question was not 
answered. The correct answer is $\theta\approx 180-44.4=135.6^\circ$.

\end{document}
